\chapter{Локальная плоскостность, ковариантное дифференцирование и кривизна}

\lectureoverview{Метрику всегда можно сделать евклидовой с точностью до первого порядка в одной выбранной точке. Кривизна начинается там, где эта свобода выбора координат кончается: в нарушении согласованности соседних локальных реперов на втором порядке. Мы разовьём это утверждение в три этапа. Сначала введём нормальные координаты, затем ковариантную производную и её связность, и наконец некоммутативность двух ковариантных производных, которая определяет тензор Римана и придаёт приливной гравитации её инвариантный смысл.}

\section{Необходимое бремя индексов}

Общая теория относительности трудна довольно особым образом. Её основные идеи не так уж непроницаемы, но вычисление быстро оказывается переполненным индексами, производными и символами. Существуют и более сжатые языки: вильбейны, дифференциальные формы, спиноры, твисторы и родственные средства позволяют упаковывать длинные покомпонентные уравнения в компактные объекты. Однако изучение этих языков само по себе составляет существенную задачу, а компактные обозначения могут скрыть именно тот механизм, который при первом знакомстве важнее всего понять.

Уравнения Максвелла дают здесь образец. Если записывать их покомпонентно, получается несколько отдельных уравнений. Векторная запись собирает их в уравнения для дивергенции и ротора, и структуру становится легче увидеть. В общей теории относительности возможны аналогичные сжатия, но покомпонентный метод остаётся надёжным путём, когда действительно нужно выполнить явное вычисление. Поэтому мы будем держать индексы на виду. Когда их грамматика станет привычной, понятийный путь окажется гораздо более гладким, хотя арифметика всё ещё может занимать целые страницы.

Символьные компьютерные пакеты могут помочь с этой арифметикой. Задав метрику \(g_{mn}(x)\), такой пакет может построить символы Кристоффеля, тензоры Римана, тензоры Риччи и тензоры Эйнштейна. Затем он может подставить эти выражения в уравнения поля. Это не устраняет необходимости понимать, что означают эти объекты. Это лишь отделяет механическое развёртывание от геометрии.

\begin{classroomqa}\audiencequestion{Каковы были активные области исследований в общей теории относительности во время этой лекции?}\lecturerresponse{В лекции было выделено несколько направлений. Квантовая гравитация и квантовая физика чёрных дыр были центральными проблемами фундаментальной теории. Классическая астрофизика и численная теория относительности изучали такие системы, как обращающиеся по орбитам и сливающиеся чёрные дыры, где решение уравнений Эйнштейна и предсказание гравитационного излучения требуют гораздо большего, чем слепо скармливать уравнения компьютеру. Те же математические структуры находили применение и в физике конденсированного состояния, и в гидродинамике. Историческая сторона была не менее примечательна: общая теория относительности перешла от сравнительно небольшой специальности в 1960-е годы к одному из центральных языков теоретической физики, тесно связанному с квантовой теорией поля.\lecturetimestamp{00:05:15}}\end{classroomqa}

Это последняя лекция, посвящённая самой по себе римановой геометрии. Начиная со следующей лекции, геометрия будет интерпретироваться как гравитация. Связь уже видна. Различение подлинного гравитационного поля и кажущегося поля, вызванного ускоренными или криволинейными координатами, математически есть та же задача, что и различение искривлённой метрики и плоской метрики, записанной в неудобных координатах.

\section{Инвариантный вопрос}

Двумерная геометрия плоска, когда евклидова геометрия верна внутренним образом: прямые, параллельные прямые, соотношения в треугольниках и измерения расстояний — всё имеет евклидову форму. Заманчиво сказать, что плоскостность означает \(g_{mn}=\delta_{mn}\), но это утверждение зависит от координат. Даже у обычной плоскости в полярных координатах метрика зависит от положения. Правильное утверждение состоит в том, что плоская область допускает координаты, в которых метрика постоянна и евклидова во всей этой области.

Предположим вместо этого, что метрика задана нам в некоторых произвольных координатах. Плохим тестом было бы перебирать системы координат одну за другой, каждый раз преобразуя метрику и надеясь в конце концов найти декартову форму. Возможностей бесконечно много, так что это не практический критерий.

Правильная стратегия — построить диагностический объект из \(g_{mn}\) и его производных. Нам нужен объект \(R\) с двумя решающими свойствами: \[
\begin{aligned}
R=0
&\quad\Longleftrightarrow\quad
\text{locally flat geometry},\\
R\neq 0 \text{ at } P
&\quad\Longrightarrow\quad
\text{curvature at } P.
\end{aligned}
\] Самое главное, \(R\) должен быть тензором. Тензор, который обращается в нуль в одной системе координат, обращается в нуль в любой системе координат, так что никакого перебора координат не требуется.

\begin{figure}[H]\centering
\includegraphics[width=\linewidth]{lecture_03_figure_01.png}\caption{Искомый инвариантный критерий: метрика не обязана явно совпадать с \(\delta_{mn}\), тогда как обращающийся в нуль объект кривизны должен указывать на локально плоскую геометрию.}\lectureframe{00:18:20}\end{figure}

Уточнение \emph{локально} существенно. Обращение тензора Римана в нуль гарантирует локально евклидову геометрию; глобальные отождествления всё ещё могут давать нетривиальную топологию. Например, плоский цилиндр или плоский тор могут быть локально евклидовыми, не будучи глобально плоскостью.

\section{Метрическая геометрия в малом}

Полезно посмотреть, что меняется с размерностью. Нульмерное пространство — это просто точка. Связное одномерное риманово пространство можно локально выпрямить в прямую; глобально оно может быть интервалом, прямой или петлёй. Жук, ограниченный замкнутой одномерной дорожкой, может измерить её длину, но изгибание этой дорожки в окружающем пространстве не меняет её внутреннюю метрику.

В двух измерениях внутренняя геометрия гораздо богаче: плоскость, сфера, неровная земная поверхность и тор могут различаться уже локально, и их глобальные топологии тоже могут быть различны. Метрика — это инструмент, который фиксирует их бесконечно малые расстояния: \begin{equation}
 ds^2=g_{mn}(x)\,dx^m dx^n.
\label{eq:l3-line-element}
\end{equation} Величина \(ds^2\) является скаляром. Следовательно, изменения \(dx^m\) при преобразовании координат в точности компенсируются преобразованием \(g_{mn}\).

\begin{figure}[H]\centering
\includegraphics[width=\linewidth]{lecture_03_figure_02.png}\caption{Инвариантный элемент длины, записанный через зависящую от положения метрику.}\lectureframe{00:19:00}\end{figure}

Метрика невырождена и потому имеет обратную: \begin{equation}
 g^{mn}g_{nr}=\delta^m{}_r.
\label{eq:l3-inverse-metric}
\end{equation} Смешанный тензор Кронекера \(\delta^m{}_r\) есть тождественное отображение. В обычной римановой геометрии метрика также положительно определена, \begin{equation}
 g_{mn}u^m u^n>0
\qquad\text{for every nonzero }u^m.
\label{eq:l3-positive}
\end{equation} Позднее лоренцево пространство-время ослабит это условие положительности, сохранив невырожденность.

\begin{figure}[H]\centering
\includegraphics[width=\linewidth]{lecture_03_figure_03.png}\caption{Обратная метрика и тензор тождества на доске.}\lectureframe{00:20:40}\end{figure}

Метрика отождествляет векторы и ковекторы: \begin{equation}
 V_m=g_{mn}V^n,
 \qquad
 V^m=g^{mn}V_n.
\label{eq:l3-raise-lower}
\end{equation} Перенос индекса, следовательно, есть операция, выполняемая метрикой, а не типографическое изменение.

\begin{figure}[H]\centering
\includegraphics[width=\linewidth]{lecture_03_figure_04.png}\caption{Поднятие и опускание векторного индекса с помощью метрики и обратной метрики.}\lectureframe{00:24:40}\end{figure}

Если все компоненты \(g_{mn}\) постоянны в некоторой области, то постоянное линейное преобразование координат может там диагонализовать и нормировать метрику. Тогда геометрия плоская. Следовательно, содержательная информация должна быть связана с тем, как метрика меняется от точки к точке.

\section{Нормальные координаты: евклидовость до первого порядка}

Выберем точку \(P\). Теорема о нормальных координатах утверждает, что можно построить координаты с центром в \(P\), для которых \begin{align}
 g_{mn}(P)&=\delta_{mn}, \label{eq:l3-normal-metric}\\
 \partial_r g_{mn}(P)&=0. \label{eq:l3-normal-first}
\end{align} В лекции они называются гауссовыми нормальными координатами. Геометрически это означает, что через \(P\) проводят оси и продолжают их координатные линии настолько прямыми, насколько позволяет геометрия. Точное понятие прямизны далее будет выражено через геодезическое движение.

\begin{figure}[H]\centering
\includegraphics[width=\linewidth]{lecture_03_figure_05.png}\caption{Нормальные координаты, выбранные так, что в выбранной точке метрика евклидова.}\lectureframe{00:28:50}\end{figure}

Такой репер не единственен. Ортогональный поворот в \(P\) сохраняет \(g_{mn}(P)=\delta_{mn}\). Важно не направление осей, а существование репера, в котором выполняются и \eqref{eq:l3-normal-metric}, и \eqref{eq:l3-normal-first}.

\begin{classroomqa}\audiencequestion{Эти условия выполняются только в выбранной точке или во всей окрестности?}\lecturerresponse{Они выполняются в выбранной точке \(P\). Нормальный репер можно построить с центром в любой точке, но, вообще говоря, нельзя использовать одну систему координат на искривлённой области и сохранять метрику евклидовой, а все первые производные равными нулю всюду. Если продолжить те же координаты в сторону от \(P\), проявятся высшие производные метрики.\lecturetimestamp{00:29:53}}\end{classroomqa}

В \(P\) вторые производные, вообще говоря, неизбежны: \begin{equation}
 \partial_r\partial_s g_{mn}(P)\neq 0
 \qquad\text{in general}.
\label{eq:l3-normal-second}
\end{equation} Таким образом, в нулевом и первом порядке в одной точке метрика не несёт инвариантной информации о кривизне; свобода выбора координат может её устранить. Первая возможная инвариантная информация появляется на втором порядке.

\begin{figure}[H]\centering
\includegraphics[width=\linewidth]{lecture_03_figure_06.png}\caption{Условия нормальных координат: метрика евклидова, её первые производные равны нулю, а вторые производные не обязаны обращаться в нуль.}\lectureframe{00:30:50}\end{figure}

\subsection{Почему координатной свободы достаточно}

Пусть \(P\) будет общим началом координат \(x^m\) и \(y^m\).
После того как с помощью линейного преобразования мы приведем метрику в начале координат к нужному виду, рассмотрим квадратичную часть дальнейшего преобразования координат:
\begin{equation}
 x^m
 =y^m+C^m{}_{nr}y^n y^r+O(y^3),
 \qquad C^m{}_{nr}=C^m{}_{rn}.
\label{eq:l3-quadratic-transform}
\end{equation}
В четырех измерениях симметричная пара \(nr\) имеет
\[
\frac{4(4+1)}{2}=10
\]
независимых вариантов. Для каждого из четырех значений \(m\) имеется
десять коэффициентов \(C^m{}_{nr}\), что дает сорок квадратичных коэффициентов.

Симметричный метрический тензор в четырех измерениях тоже имеет десять независимых компонент. У каждой компоненты имеются четыре первые производные, так что условия \(\partial_r g_{mn}(P)=0\) составляют сорок уравнений. Совпадение числа параметров показывает, что квадратичная координатная свобода имеет в точности нужный размер, чтобы наложить условия на первые производные. Теорема о нормальных координатах дает утверждение о существовании; один лишь подсчет служит полезной структурной проверкой, а не полным доказательством разрешимости произвольной нелинейной системы.

\begin{figure}[H]
\centering
\includegraphics[width=\linewidth]{lecture_03_figure_07.png}
\caption{Квадратичное преобразование координат и десять независимых квадратичных мономов в четырех измерениях.}
\lectureframe{00:33:50}
\end{figure}

\begin{classroomqa}
\audiencequestion{Можно ли рассматривать \(x^m=y^m+C^m{}_{nr}y^n y^r+\cdots\) как преобразование координат?}
\lecturerresponse{Да. Коэффициенты \(C^m{}_{nr}\) выбираются так, чтобы в новых координатах первые производные метрики обращались в нуль в общем начале координат. Линейная часть сохраняет якобиан обратимым вблизи этой точки, а квадратичная часть подстраивает первые производные метрики. \lecturetimestamp{00:35:41}}
\end{classroomqa}

\begin{classroomqa}
\audiencequestion{Именно это и имеется в виду, когда говорят, что гладкая геометрия локально плоска?}
\lecturerresponse{Это и есть существенное утверждение первого порядка. В достаточно малом масштабе гладкая поверхность аппроксимируется своей касательной плоскостью, и нормальные координаты выражают этот факт через \(g_{mn}(P)=\delta_{mn}\) и \(\partial_r g_{mn}(P)=0\). Точная плоскостность в некоторой окрестности — более сильное условие: для этого также требуется, чтобы тензор кривизны там обращался в нуль.
\lecturetimestamp{00:36:19}}
\end{classroomqa}

\begin{classroomqa}
\audiencequestion{Действительно ли равенство между сорока коэффициентами и сорока уравнениями гарантирует, что уравнения можно решить?}
\lecturerresponse{Подсчет выявляет требуемую координатную свободу, а теорема гарантирует локальное построение для гладкой невырожденной метрики. В лекции предлагается в качестве упражнения проверить двумерный случай непосредственно. Рассуждение не зависит от четырех измерений; в \(d\) измерениях та же симметрия и та же координатная свобода согласуются друг с другом. \lecturetimestamp{00:36:45}}
\end{classroomqa}

\begin{classroomqa}
\audiencequestion{Почему в разложении нет постоянного члена, и какой член является членом первого порядка?}
\lecturerresponse{Член \(y^m\) является членом первого порядка. Постоянный член был опущен, потому что начала обеих систем координат намеренно выбраны так, чтобы представлять одну и ту же точку, так что \(y=0\) влечет \(x=0\). Этот выбор не ограничивает общности. \lecturetimestamp{00:38:52}}
\end{classroomqa}

\section{Почему покомпонентное дифференцирование не работает}

Рассмотрим теперь ковариантное векторное поле \(V_m(x)\). В прямолинейных декартовых координатах на плоском пространстве геометрически постоянный вектор имеет постоянные компоненты. Поэтому можно было бы попытаться определить его производную как \(\partial_rV_m\). Трудность возникает, как только координатная сетка искривляется.

\begin{figure}[H]
\centering
\includegraphics[width=\linewidth]{lecture_03_figure_08.png}
\caption{Криволинейная координатная сетка, наложенная на локально прямолинейный репер.}
\lectureframe{00:40:10}
\end{figure}

Рассмотрим один и тот же геометрический вектор в соседних точках. Относительно фиксированного декартова репера его компоненты не меняются. Однако в криволинейном базисе сами базисные векторы поворачиваются при перемещении, так что набор компонент меняется, хотя сам вектор — нет.

\begin{figure}[H]
\centering
\includegraphics[width=\linewidth]{lecture_03_figure_09.png}
\caption{Один и тот же вектор в соседних точках приобретает разные компоненты, потому что локальный координатный базис повернулся.}
\lectureframe{00:42:30}
\end{figure}

Наивное условие
\begin{equation}
 \partial_rV_m=0
\label{eq:l3-naive-derivative}
\end{equation}
поэтому может выполняться в одной системе координат и не выполняться в другой. Если бы \(\partial_rV_m\) было тензором второго ранга, это было бы невозможно: обращение в нуль в одном репере влекло бы обращение в нуль в любом репере. Обычные частные производные компонент ковариантного вектора поэтому не являются компонентами тензора.

\begin{figure}[H]
\centering
\includegraphics[width=\linewidth]{lecture_03_figure_10.png}
\caption{Намеренно наивное условие нулевой производной рядом с изменяющимся криволинейным базисом.}
\lectureframe{00:44:00}
\end{figure}

Нам нужна производная, которая устраняет кажущееся изменение, вызванное базисом, сохраняя при этом подлинное изменение векторного поля.

\section{Определение ковариантной производной}

Дифференцирование в \(P\) сравнивает значения только в сколь угодно малой окрестности точки \(P\). Поэтому можно начать с нормальных координат, центрированных в ней. Возьмем вектор в соседней точке, перенесем его начало обратно в \(P\) и во время переноса будем держать его компоненты в нормальных координатах фиксированными. Сравнение перенесенного вектора с тем, который уже задан в \(P\), а затем деление на расстояние между точками дает обычную производную в этом локально прямолинейном репере.

\begin{figure}[H]
\centering
\includegraphics[width=\linewidth]{lecture_03_figure_11.png}
\caption{Локальное сравнение: соседний вектор сначала переносится в нормальных координатах, и лишь затем берется разность.}
\lectureframe{00:46:50}
\end{figure}

Теперь выразим эту геометрически определенную производную в произвольных координатах. Имеются два вклада: компоненты \(V_m\) могут изменяться, и координатный базис тоже может изменяться. Структура индексов определяет вид поправки:
\begin{equation}
 \nabla_rV_m
 =\partial_rV_m-\Gamma^t{}_{rm}V_t.
\label{eq:l3-covector-derivative}
\end{equation}
Величины \(\Gamma^t{}_{rm}\) — это коэффициенты связности, также называемые символами Кристоффеля. Знак минус соответствует нижнему индексу. В центре нормальных координат, в точке \(P\),
\begin{equation}
 \Gamma^t{}_{rm}(P)=0,
\qquad
 \nabla_rV_m(P)=\partial_rV_m(P).
\label{eq:l3-gamma-normal-zero}
\end{equation}

\begin{figure}[H]
\centering
\includegraphics[width=\linewidth]{lecture_03_figure_12.png}
\caption{Ковариантная производная ковектора: обычное изменение компонент минус изменение, индуцированное движущимся базисом.}
\lectureframe{00:49:50}
\end{figure}

\begin{classroomqa}
\audiencequestion{Верна ли формула \(\nabla_rV_m=\partial_rV_m-\Gamma^t{}_{rm}V_t\) в произвольных координатах и во что она превращается в нормальных координатах?}
\lecturerresponse{Полная формула верна в произвольных координатах. В выбранной точке в нормальных координатах первые производные метрики, а значит и символы Кристоффеля, обращаются в нуль, так что остается только \(\partial_rV_m\). Вне этой точки даже в той же карте нормальных координат коэффициенты связности, вообще говоря, ненулевые.
\lecturetimestamp{00:50:22}}
\end{classroomqa}

Ковариантная производная так называется потому, что весь объект \(\nabla_rV_m\) преобразуется ковариантно как тензор; это название не является просто замечанием о нижнем индексе производной. Член с символами Кристоффеля как раз и исправляет нетензорное преобразование частной производной.

\begin{figure}[H]
\centering
\includegraphics[width=\linewidth]{lecture_03_figure_13.png}
\caption{Полная производная ковектора и соответствующая поправка связности.}
\lectureframe{00:51:30}
\end{figure}

\begin{classroomqa}
\audiencequestion{Связано ли это с конвективной производной, используемой в течении жидкости?}
\lecturerresponse{В лекции на это отвечают осторожно: связь есть. В механике жидкости также различают изменение, переносимое течением, и изменение, наблюдаемое в фиксированной координатной точке, а ковариантные производные возникают в уравнениях жидкости в криволинейных координатах или на искривленных пространствах. Эти операторы по определению не тождественны, но оба исправляют наивную производную с учетом структуры, движущейся вместе с описанием.
\lecturetimestamp{00:52:31}}
\end{classroomqa}

Для используемой здесь связи без кручения,
\begin{equation}
 \Gamma^t{}_{rm}=\Gamma^t{}_{mr}.
\label{eq:l3-torsion-free}
\end{equation}
Более общие геометрии могут обладать кручением и не обязаны удовлетворять этой симметрии. Рассматриваемая здесь связь — именно та, которая относится к стандартной метрической формулировке общей теории относительности.

\begin{figure}[H]
\centering
\includegraphics[width=\linewidth]{lecture_03_figure_14.png}
\caption{Символы Кристоффеля как симметричные коэффициенты связности геометрии без кручения.}
\lectureframe{00:54:30}
\end{figure}

Есть полезная физическая аналогия. Нормальные координаты играют роль свободно падающей системы отсчета в одном событии: там координатную поправку можно обратить в нуль. Произвольные криволинейные координаты напоминают ускоренный лифт, где появляются кажущиеся силы, потому что изменяется сама система отсчета. Эта аналогия локальна, и кривизна покажет, почему одна свободно падающая система отсчета не может устранить гравитацию во всей протяженной области.

\subsection{По одной поправке на каждый тензорный индекс}

Для ковариантного тензора ранга два \(T_{mn}\) каждый нижний индекс получает свою собственную поправку связности:
\begin{equation}
 \nabla_sT_{mn}
 =\partial_sT_{mn}
 -\Gamma^t{}_{sm}T_{tn}
 -\Gamma^t{}_{sn}T_{mt}.
\label{eq:l3-rank-two-derivative}
\end{equation}
Чтобы запомнить правило, держите один индекс фиксированным как пассивный, пока дифференцируете по другому, а затем повторите это для второго индекса. Контравариантные индексы получают соответствующие члены с противоположным знаком.

\begin{figure}[H]
\centering
\includegraphics[width=\linewidth]{lecture_03_figure_15.png}
\caption{Ковариантная производная ковариантного тензора ранга два: по одной поправке Кристоффеля на каждый нижний индекс.}
\lectureframe{01:01:50}
\end{figure}

\begin{classroomqa}
\audiencequestion{Относится ли сам символ \(\partial_r\) только к нормальным координатам, и ограничена ли тензорная формула этой системой отсчета?}
\lecturerresponse{Нет. Обозначение \(\partial_r=\partial/\partial x^r\) имеет смысл в любой системе координат, и полная формула ковариантной производной записывается для произвольных координат. Только в выбранной точке нормального репера поправки Кристоффеля исчезают. \lecturetimestamp{00:57:22}}
\end{classroomqa}

\begin{classroomqa}
\audiencequestion{Когда в нормальных координатах ковариантная производная сводится к частной производной, происходит ли это сведение только в выбранной точке? Являются ли символы Кристоффеля тензорами?}
\lecturerresponse{Да, это сведение имеет место поточечно. Символы Кристоффеля не являются тензорами: они могут обращаться в нуль в \(P\) в нормальных координатах и быть ненулевыми в той же точке в другой системе координат. Именно их нетензорное преобразование и позволяет комбинации в \eqref{eq:l3-covector-derivative} преобразовываться как тензор.
\lecturetimestamp{00:58:13}}
\end{classroomqa}

\begin{classroomqa}
\audiencequestion{Для чего на практике будет использоваться ковариантная производная?}
\lecturerresponse{Для уравнений поля. Гравитационные поля меняются от точки к точке, поэтому их уравнения должны содержать производные. В то же время эти уравнения должны иметь одно и то же содержание в любой системе координат. Ковариантное дифференцирование переводит тензоры в тензоры и тем самым допускает дифференциальные уравнения, истинность которых не зависит от какой-либо выделенной системы отсчёта. \lecturetimestamp{01:02:35}}
\end{classroomqa}

\begin{classroomqa}
\audiencequestion{Можно ли умножить ковариантную производную на координатный дифференциал, чтобы описать изменение при переходе между соседними точками?}
\lecturerresponse{Да. При бесконечно малом смещении \(dx^s\) ковариантное изменение равно \(dx^s\nabla_sT\). Частный член измеряет изменение компонент, тогда как члены со связностью учитывают изменение координатного базиса. Даже в плоском пространстве, записанном в полярных координатах, это различие необходимо. \lecturetimestamp{01:03:22}}
\end{classroomqa}

\section{Метрическая совместимость определяет связность}

Применим к метрике правило для тензора ранга два. В \(P\) в нормальных координатах \(\partial_sg_{mn}=0\) и \(\Gamma^t{}_{sm}=0\), так что \(\nabla_sg_{mn}=0\). Поскольку ковариантная производная является тензором, этот поточечный результат справедлив в любой системе координат:
\begin{equation}
 \nabla_sg_{mn}=0.
\label{eq:l3-metric-compatibility}
\end{equation}
Это и есть метрическая совместимость. Она означает, что параллельный перенос сохраняет метрику и, следовательно, сохраняет длины и углы.

\begin{figure}[H]
\centering
\includegraphics[width=\linewidth]{lecture_03_figure_16.png}
\caption{Уравнение метрической совместимости служит отправной точкой для вывода связности из метрики.}
\lectureframe{01:06:20}
\end{figure}

Раскрывая \eqref{eq:l3-metric-compatibility}, получаем
\begin{equation}
 0
 =\partial_sg_{mn}
 -\Gamma^t{}_{sm}g_{tn}
 -\Gamma^t{}_{sn}g_{mt}.
\label{eq:l3-compat-expanded}
\end{equation}
Запишем ещё две копии с циклической перестановкой индексов производной и метрики:
\begin{align}
 0
 &=\partial_mg_{sn}
 -\Gamma^t{}_{ms}g_{tn}
 -\Gamma^t{}_{mn}g_{st},
\label{eq:l3-compat-m}\\
 0
 &=\partial_ng_{sm}
 -\Gamma^t{}_{ns}g_{tm}
 -\Gamma^t{}_{nm}g_{st}.
\label{eq:l3-compat-n}
\end{align}

\begin{figure}[H]
\centering
\includegraphics[width=\linewidth]{lecture_03_figure_17.png}
\caption{Три уравнения метрической совместимости с переставленными индексами.}
\lectureframe{01:10:20}
\end{figure}

Сложим \eqref{eq:l3-compat-m} и \eqref{eq:l3-compat-n}, затем вычтем \eqref{eq:l3-compat-expanded}. Используя симметрию при отсутствии кручения \eqref{eq:l3-torsion-free}, нежелательные члены со связностью сокращаются:
\begin{equation}
 \partial_mg_{sn}+\partial_ng_{sm}-\partial_sg_{mn}
 =2\Gamma^t{}_{mn}g_{st}.
\label{eq:l3-isolate-gamma-lowered}
\end{equation}

\begin{figure}[H]
\centering
\includegraphics[width=\linewidth]{lecture_03_figure_18.png}
\caption{При сложении двух уравнений совместимости и вычитании третьего нежелательные члены с символами Кристоффеля сокращаются.}
\lectureframe{01:12:10}
\end{figure}

Умножение на \(g^{rs}\) позволяет выделить связность:
\begin{equation}
 \boxed{
 \Gamma^r{}_{mn}
 =\frac12 g^{rs}
 \left(
 \partial_mg_{sn}
 +\partial_ng_{sm}
 -\partial_sg_{mn}
 \right)}
\label{eq:l3-christoffel}
\end{equation}
Итак, связность Леви-Чивиты строится из метрики и её первых производных. Это и есть обещанное объяснение того, почему символы Кристоффеля обращаются в нуль в центре нормальных координат.

\begin{figure}[H]
\centering
\includegraphics[width=\linewidth]{lecture_03_figure_19.png}
\caption{Обратная метрика выделяет символ Кристоффеля из уравнения для связности с пониженным индексом.}
\lectureframe{01:16:10}
\end{figure}

\begin{classroomqa}
\audiencequestion{Почему все три уравнения метрической совместимости были приравнены к нулю?}
\lecturerresponse{В выбранной точке в нормальных координатах обычные производные метрики обращаются в нуль, и символы Кристоффеля тоже. Следовательно, ковариантная производная метрики там равна нулю. Чтобы установить это утверждение в любой другой точке, выберем нормальный репер, центрированный в этой точке, и повторим рассуждение. Поскольку \(\nabla_sg_{mn}\) — тензор, его обращение в нуль не зависит от системы координат. \lecturetimestamp{01:14:27}}
\end{classroomqa}

\begin{classroomqa}
\audiencequestion{Каковы размерности метрики и массива символов Кристоффеля, и почему в четырёх измерениях имеется сорок независимых компонент Кристоффеля?}
\lecturerresponse{В \(d\) измерениях \(g_{mn}\) — симметричная матрица \(d\times d\). У связности три индекса. Для каждого фиксированного верхнего индекса \(r\) симметричная нижняя пара \(mn\) имеет \(d(d+1)/2\) вариантов. Поэтому связность без кручения имеет \(d^2(d+1)/2\) независимых компонент, что даёт сорок при \(d=4\).
\lecturetimestamp{01:16:18}}
\end{classroomqa}

\begin{figure}[H]
\centering
\includegraphics[width=\linewidth]{lecture_03_figure_20.png}
\caption{Запись на доске при обсуждении числа компонент: десять вариантов для симметричной нижней пары и четыре варианта для верхнего индекса.}
\lectureframe{01:16:30}
\end{figure}

\begin{classroomqa}
\audiencequestion{Какая метрика фигурирует в
\eqref{eq:l3-christoffel}? Нужно ли в каждой точке использовать другую метрику в нормальных координатах?}
\lecturerresponse{В формуле используются компоненты метрики в тех глобальных или локальных координатах, которые выбраны для вычисления. Нормальные координаты — лишь локальный приём, используемый для вывода правила: перейти к нормальному реперу в одной точке, продифференцировать там, преобразовать результат обратно, а затем отбросить вспомогательный репер. После получения \eqref{eq:l3-christoffel} его используют непосредственно в произвольных координатах, представляющих интерес. \lecturetimestamp{01:17:21}}
\end{classroomqa}

\begin{classroomqa}
\audiencequestion{Следует ли из \(\nabla_sg_{mn}=0\) или из формулы для символов Кристоффеля, что пространство плоское?}
\lecturerresponse{Нет. Метрическая совместимость выполняется для любой связности Леви-Чивиты, как в искривлённом, так и в плоском случае. Символы Кристоффеля могут быть ненулевыми даже на плоскости, записанной в полярных координатах. Если все \(\Gamma^r{}_{mn}\) обращаются в нуль во всей области в одной координатной карте, то эта карта там является декартовой и область плоская; но сама по себе их ненулевость не позволяет диагностировать кривизну. По-прежнему нужен тензор, построенный из производных и произведений \(\Gamma\). \lecturetimestamp{01:19:27}}
\end{classroomqa}

\section{Конус выявляет то, что связность выявить не может}

Рассмотрим сглаженный конус. Вдали от его вершины его можно аппроксимировать обычным бумажным конусом. Вырежем сектор из плоского диска и склеим образовавшиеся края. Оставшаяся бумага нигде не была растянута; локально каждая точка вдали от шва и вершины по-прежнему плоская. Отсутствующий сектор характеризуется углом дефицита \(\Delta\).

\begin{figure}[H]
\centering
\includegraphics[width=\linewidth]{lecture_03_figure_21.png}
\caption{Конус вводится как геометрия, локально плоская вдали от вершины, но имеющая нетривиальную кривизну, сосредоточенную в вершине.}
\lectureframe{01:23:00}
\end{figure}

\begin{figure}[H]
\centering
\includegraphics[width=\linewidth]{lecture_03_figure_22.png}
\caption{Вырезание сектора из диска и отождествление его краёв порождают угол дефицита конуса.}
\lectureframe{01:24:20}
\end{figure}

Бумажная конструкция делает это различие наглядным. На разрезанном листе перенесите вектор, сохраняя его параллельным в обычной плоскости. После отождествления краёв сектора путь, обходящий вершину, возвращается в исходную точку, причём перенесённый вектор оказывается повёрнутым относительно своего первоначального направления.

\begin{figure}[H]
\centering
\includegraphics[width=0.78\linewidth]{lecture_03_figure_23.png}
\caption{Демонстрация бумажного конуса: плоский лист с вырезанным сектором становится конусом, когда его разрезанные края соединяют.}
\lectureframe{01:25:10}
\end{figure}

\begin{figure}[H]
\centering
\includegraphics[width=\linewidth]{lecture_03_figure_24.png}
\caption{Векторы, нарисованные на вырезанном секторе до отождествления краёв.}
\lectureframe{01:27:40}
\end{figure}

Этот поворот — голономия. Он не вызван неудачным выбором координат на одном из участков. Каждый малый участок можно рассматривать в локально плоском репере. Эффект возникает только после композиции этих локальных сравнений вокруг замкнутого контура. Для идеального конуса суммарный поворот равен углу дефицита с точностью до ориентации.

\begin{figure}[H]
\centering
\includegraphics[width=\linewidth]{lecture_03_figure_25.png}
\caption{Параллельный перенос вокруг вершины конуса возвращает вектор с повёрнутым направлением: голономия фиксирует заключённую внутри кривизну.}
\lectureframe{01:29:10}
\end{figure}

\section{Два пути вокруг бесконечно малого контура}

Заменим конечный контур крошечным координатным прямоугольником. Начиная из одного угла, перенесём сначала в направлении \(r\), а затем в направлении \(s\). Сравним этот результат с переносом сначала в направлении \(s\), а затем в направлении \(r\). В плоских декартовых координатах оба пути дают один и тот же результат. В искривлённой геометрии, вообще говоря, это не так.

\begin{figure}[H]
\centering
\includegraphics[width=\linewidth]{lecture_03_figure_26.png}
\caption{Два порядка бесконечно малого переноса вокруг координатного прямоугольника.}
\lectureframe{01:31:40}
\end{figure}

Алгебраической мерой служит коммутатор ковариантных производных:
\begin{equation}
 \left(\nabla_s\nabla_r-\nabla_r\nabla_s\right)V_n.
\label{eq:l3-commutator}
\end{equation}
Обычные смешанные частные производные коммутируют при действии на гладкие функции. Здесь нарушение возникает из-за того, что дифференцируются не только компоненты вектора, но и члены со связностью.

\begin{figure}[H]
\centering
\includegraphics[width=\linewidth]{lecture_03_figure_27.png}
\caption{Бесконечно малый контур, представленный разностью двух порядков ковариантного дифференцирования.}
\lectureframe{01:33:20}
\end{figure}

Чтобы вычислить его, сначала рассмотрим
\[
 T_{rn}=\nabla_rV_n
\]
как ковариантный тензор ранга два. Затем используем \eqref{eq:l3-rank-two-derivative} для дифференцирования \(T_{rn}\) по направлению \(s\). Повторим это, поменяв местами \(r\) и \(s\), и вычтем. Вычисление длинное, но механическое: нужно раскрыть выражения, применить правило произведения и сократить члены, содержащие вторые частные производные \(V_n\).

\begin{figure}[H]
\centering
\includegraphics[width=\linewidth]{lecture_03_figure_28.png}
\caption{Начало вычисления двойной ковариантной производной, где \(\nabla_rV_n\) рассматривается как тензор ранга два.}
\lectureframe{01:35:40}
\end{figure}

Используя принятое на доске соглашение, получаем
\begin{equation}
 \left(\nabla_s\nabla_r-\nabla_r\nabla_s\right)V_n
 =R_{srn}{}^tV_t.
\label{eq:l3-riemann-commutator}
\end{equation}

\begin{figure}[H]
\centering
\includegraphics[width=\linewidth]{lecture_03_figure_29.png}
\caption{Коммутатор ковариантных производных определяет действие тензора кривизны на ковектор.}
\lectureframe{01:38:30}
\end{figure}

Для такого порядка в коммутаторе,
\begin{equation}
\begin{split}
 R_{srn}{}^t
={}&
 \partial_r\Gamma^t{}_{sn}
 -\partial_s\Gamma^t{}_{rn}\\
&+\Gamma^p{}_{sn}\Gamma^t{}_{pr}
 -\Gamma^p{}_{rn}\Gamma^t{}_{ps}.
\end{split}
\label{eq:l3-riemann}
\end{equation}
В других текстах знак может быть противоположным или индексы могут быть расположены иначе; соглашение фиксируется тем, что явно записываются и \eqref{eq:l3-riemann-commutator}, и \eqref{eq:l3-riemann}.

\begin{figure}[H]
\centering
\includegraphics[width=\linewidth]{lecture_03_figure_30.png}
\caption{Тензор Римана: две производные связности и два квадратичных по связности члена.}
\lectureframe{01:39:40}
\end{figure}

Уравнение \eqref{eq:l3-christoffel} показывает, почему кривизна впервые возникает лишь во втором порядке по метрике. Каждый \(\Gamma\) содержит первую производную от \(g\). Поэтому производные члены в \eqref{eq:l3-riemann} содержат вторые производные \(g\), тогда как квадратичные члены содержат произведения первых производных. В нормальных координатах в \(P\) сама связность обращается в нуль, но её производные, вообще говоря, не обращаются, так что кривизна не исчезает:
\[
\begin{aligned}
 \Gamma^t{}_{mn}(P)&=0,\\
 R_{srn}{}^t(P)
 &=\partial_r\Gamma^t{}_{sn}(P)
 -\partial_s\Gamma^t{}_{rn}(P).
\end{aligned}
\]
Важно не путать порядок производных: тензор Римана использует первые производные \(\Gamma\), которые относятся ко второму порядку по метрике.

Это и есть искомый критерий. Если в некоторой области каждая компонента тензора Римана обращается в нуль, эта область локально плоская. Если в какой-либо точке хотя бы одна компонента отлична от нуля, никакая замена координат не может устранить там кривизну. Система компьютерной алгебры может проверить эти компоненты для аналитически заданной метрики, но сама идея проста: параллельный перенос вокруг бесконечно малого контура не замыкается.

\section{Кривизна как приливная гравитация}

Геометрический критерий имеет прямой физический смысл. Представьте себе гибкий объект, составленный из соединённых треугольников. В плоской области его можно перемещать без деформационных напряжений. Если вдвинуть его в область, где искривлена внутренняя метрика, разные части этой конструкции уже не смогут одновременно сохранять свои предпочтительные взаимные расстояния и углы. Объекту придётся растянуться, сжаться или испытать сдвиговую деформацию.

\begin{figure}[H]
\centering
\includegraphics[width=\linewidth]{lecture_03_figure_31.png}
\caption{Связанная конструкция, перемещаемая из плоской области в искривлённую выпуклую область, должна деформироваться, чтобы следовать локальной метрике.}
\lectureframe{01:44:10}
\end{figure}

Эта деформация — геометрический аналог приливной силы. В гравитационном поле соседние свободно падающие частицы могут ускоряться друг относительно друга: в одном направлении может происходить растяжение, тогда как в другом — сжатие. Тензор Римана определяет это относительное ускорение. В терминах, которые позже будут введены как геодезическое отклонение, для касательного вектора \(u^\mu\) и вектора разделения \(\xi^\mu\),
\begin{equation}
 \frac{D^2\xi^\mu}{D\tau^2}
 =-R^\mu{}_{\nu\rho\sigma}
 u^\nu\xi^\rho u^\sigma,
\label{eq:l3-geodesic-deviation-preview}
\end{equation}
где общий знак определяется выбранным соглашением о кривизне. Это уравнение здесь не выводится, но оно точно формулирует утверждение лекции о том, что кривизна и приливное напряжение — буквально одно и то же локальное физическое явление.

Идеальное совершенно однородное гравитационное поле не имеет приливной неоднородности и может быть локально устранено переходом в равномерно ускоренную систему отсчёта. Реальные поля, создаваемые изолированными массами, не являются совершенно однородными. Поле Земли направлено приблизительно радиально и меняется от точки к точке, поэтому соседние траектории свободного падения сходятся или расходятся.

\begin{figure}[H]
\centering
\includegraphics[width=\linewidth]{lecture_03_figure_32.png}
\caption{Радиальное гравитационное поле меняется на протяжении протяжённого тела, порождая относительное ускорение, отождествляемое с приливной гравитацией.}
\lectureframe{01:48:40}
\end{figure}

\begin{classroomqa}
\audiencequestion{Что отличает подлинное гравитационное поле от кажущегося поля, видимого в ускоренной системе отсчёта?}
\lecturerresponse{Равномерно ускоренная система отсчёта может имитировать однородное гравитационное поле в локальных экспериментах, но она не порождает приливного относительного ускорения. Поле, создаваемое реальным веществом, меняется от места к месту и потому создаёт приливные силы. Эти силы служат инвариантным свидетельством того, что эффект не является всего лишь преобразованием координат. Совершенно однородное реальное поле — идеализация; вблизи достаточно малого участка поверхности большого тела действительное поле может быть лишь приблизительно однородным. \lecturetimestamp{01:46:54}}
\end{classroomqa}

\begin{classroomqa}
\audiencequestion{Почему для приливных эффектов важен размер падающего объекта?}
\lecturerresponse{Кривизна связана со сравнением соседних точек. Объект, который намного меньше масштаба длины, на котором меняется поле, испытывает почти одно и то же ускорение во всех своих частях и мало деформируется. Протяжённый объект испытывает большую разность ускорений. Свободно падающая бактерия рядом с Землёй испытывает пренебрежимо малое приливное напряжение, тогда как воображаемый в лекции человек длиной в две тысячи миль был бы сильно деформирован. Кривизна имеет локальное значение, но накопленная деформация растёт с расстоянием между сравниваемыми точками.
\lecturetimestamp{01:49:01}}
\end{classroomqa}

\begin{classroomqa}
\audiencequestion{Если для проверки плоскостности нужно проверять кривизну в каждой точке, не остаётся ли критерий по-прежнему бесконечной задачей?}
\lecturerresponse{Для метрики, заданной аналитически, производные и компоненты кривизны являются аналитическими выражениями. Их вычисляют один раз как функции положения, и часто можно увидеть, обращаются ли они тождественно в нуль. Если метрика задана только в виде численной таблицы, тогда нужно оценивать и исследовать кривизну во всей области, покрытой выборкой, что действительно может быть затратным. Даже в этом случае задача определена гораздо лучше, чем перебор всех возможных систем координат в каждой точке.
\lecturetimestamp{01:50:37}}
\end{classroomqa}

\section{Что было установлено}

Теперь этот аргумент можно прочитать как одну непрерывную цепочку:
\begin{enumerate}
 \item В любой выбранной точке можно ввести координаты, в которых \(g_{mn}=\delta_{mn}\) и \(\partial_rg_{mn}=0\).
 \item Обычные производные компонент тензора не преобразуются как тензоры, потому что координатные базисы меняются от точки к точке.
 \item Связность компенсирует это изменение базиса и определяет ковариантное дифференцирование.
 \item Метрическая совместимость и нулевое кручение однозначно определяют связность по метрике.
 \item Две ковариантные производные не коммутируют; их коммутатор задаётся тензором Римана. Это некоммутирование измеряет голономию, вариацию метрики второго порядка и приливную гравитацию.
\end{enumerate}

Таким образом, обещанный критерий кривизны построен без выбора привилегированной глобальной системы координат. Следующий шаг — перейти от римановой геометрии с положительно определённой метрикой к лоренцеву пространству-времени, применить эти инструменты к геометрии Шварцшильда звезды или чёрной дыры и вернуться к геодезическим, рассмотрение которых было отложено в конце этой лекции.
